\scanpage{30}
les choses de façon plus intrinsèque. Considérons
l'homomorphisme canonique
\[
\unsure{\Pi|\pi}_\infty\longrightarrow\Pi_2\tag*{(16)}
\]
et soit
\[
\Pi_2^*=\Pi_2\setminus\{1\}\simeq\deleted{\unsure{(\mathbb Z/2\mathbb Z)\times\mathbb Z/2\mathbb Z}}
\inserted{(\mathbb F_2^3/\mathbb F_2)\setminus\{0\}}\tag*{(17)}
\]
c'est isom à l'ens $I=\{0,1,\infty\}$ de façon compatible
avec l'action de $\unsure{\mathfrak S|\mathscr S}_3$ [NB $\Gamma$ opère \unsure{trivialt}
sur $\Pi_2^*$ comme sur $\inserted{I=}\{0,1,\infty\}$]

Ceci dit, \deleted{\unsure{on a, si}}
\deleted{\[
 I=\{0,1,\infty\},\quad E(I)=\overline C_\infty(\overline{\mathbb Q})|I
 =\coprod_{\zeta\in I}\overline C_{\infty\,\zeta}
 =E_0\amalg E_1\amalg E_\infty\tag*{(18)}
\]
posant $\widehat\mu_\infty=\mu_\infty\amalg\{0\}$
\[
 E(I)=\{(\zeta_0,\zeta_1,\zeta_\infty)\in\widetilde\mu_\infty^3\mid
 \text{un des }\zeta_i\text{ est nul}\}\tag*{(19)}
\]}
considérons la suite \deleted{\unsure{exacte plus}} \inserted{exacte de $(\mathfrak S_3\times\Gamma)$-modules}
\[
\begin{aligned}
\deleted{\unsure{1}}0\longrightarrow\Pi_\infty\xrightarrow{2}\Pi_\infty
\xrightarrow{p}{}&\Pi_2\longrightarrow{}\\
&\rotatebox{90}{$\subset$}\\
&I=\{0,1,\infty\}
\end{aligned}\tag*{(18)}
\]
et pour $i\in I$, soit \deleted{$\unsure{Z_i\subset}$}
\[
\begin{aligned}
Z_i&\subset\Pi_\infty\\
\rotatebox{90}{$\simeq$}&\\[-3pt]
T&
\end{aligned}\tag*{(19)}
\]
% slanted
\margin{Autre façon : on utilise les trois points rat/$\mathbb Q$ de $\overline C_\infty$, \deleted{\unsure{d'origine|d'ailleurs}} \unsure{des|de} $E_i$ $(i=0,1,\infty)$ d'où $E_i\simeq\Pi_\infty/Z_i\simeq T\mathbin{\wedge}_{(\pm1)}(I\setminus\{i\})$.}
le \uline{sous-groupe de ramification} correspondant,
image dans $\Pi_\infty\simeq T^3/T$ du \deleted{\unsure{ième|i-ème}} facteur $T$
d'indice $i$ dans $T^3$. \deleted{Soit enfin} On a \unsure{donc|des|alors}
\[
E_i\simeq p^{-1}(\{i\})/Z_i\tag*{(20)}
\]
% archive sheet number: 29
